% ============================================================
% Section 2: Related Work
% The Information Budget of Graph Neural Networks
% ============================================================

\section{Related Work}
\label{sec:related}

Our Information Budget framework relates to four research areas: GNN aggregation mechanisms, homophily and heterophily analysis, feature-structure relationships, and GNN expressiveness theory.

\subsection{Graph Neural Networks and Aggregation}

\textbf{Spectral methods}: GCN~\cite{kipf2017gcn} introduced efficient graph convolution via symmetric normalization. SGC~\cite{wu2019sgc} simplified this by removing nonlinearities. These methods use \textit{replace} aggregation, where the node's own features are overwritten by aggregated neighbors.

\textbf{Spatial methods}: GraphSAGE~\cite{hamilton2017graphsage} uses neighborhood sampling with \textit{concatenation}---preserving the node's own features alongside aggregated information. GIN~\cite{xu2019gin} proved theoretical expressiveness bounds. Our ADR analysis (Section~\ref{sec:adr}) shows that the replace-vs-concatenate distinction is critical for robustness.

\textbf{Attention mechanisms}: GAT~\cite{velivckovic2018gat} learns neighbor importance. GATv2~\cite{brody2022gatv2} improved dynamic attention. However, attention cannot prevent damage when neighbors are fundamentally misleading---a problem our ADR quantifies.

\textbf{Scalability}: GraphSAINT~\cite{zeng2020graphsaint} and Cluster-GCN~\cite{chiang2019clustergcn} enable million-node training. Our OGB experiments (Section~\ref{sec:budget}) validate the Information Budget at this scale.

\subsection{Homophily and Heterophily}

\textbf{The homophily assumption}: Classic GNNs assume connected nodes share labels. Edge homophily $h = \frac{|\{(u,v) \in E : y_u = y_v\}|}{|E|}$ measures this tendency~\cite{mcpherson2001homophily}.

\textbf{Heterophily-aware methods}: When $h$ is low, standard GNN performance is inconsistent. H2GCN~\cite{zhu2020beyond} uses 2-hop neighbors and ego-neighbor separation. LINKX~\cite{luan2022linkx} separates features from structure entirely. MixHop~\cite{abu2019mixhop} and GPR-GNN~\cite{chien2021gprgnn} use multi-hop aggregation with learnable coefficients. FAGCN~\cite{fagcn2021} introduces frequency-adaptive filtering to handle both homophily and heterophily. GloGNN~\cite{li2022glognn} discovers global homophily patterns even in locally heterophilous graphs. ACM-GCN~\cite{acmgcn2024} uses adaptive cross-hop message passing. Our recovery-ratio analysis (Section~\ref{sec:heterophily}) diagnoses whether label homophily increases at two hops without claiming that this statistic alone selects the optimal model.

\textbf{Beyond binary homophily}: Recent work challenges the simple high/low dichotomy. Platonov et al.~\cite{platonov2023heterophily} show that homophily alone is insufficient for predicting GNN success. Ma et al.~\cite{ma2021homophily} introduce class-conditional homophily. Our Information Budget provides a complementary perspective: feature sufficiency determines the ceiling, regardless of $h$.

\subsection{Feature-Structure Relationships}

\textbf{When features are enough}: Huang et al.~\cite{huang2021mlp} showed MLPs can match GNNs on some benchmarks. GAMLP~\cite{zhang2021gamlp} precomputes propagation then trains MLP. Our work formalizes this: when $\text{Acc}_{\text{MLP}}$ is high, there is insufficient \textbf{budget} for GNN improvement.

\textbf{Feature quality metrics}: Node feature informativeness has been studied~\cite{shchur2018pitfalls}, but primarily as a data quality issue. We reconceptualize it as the \textbf{Information Budget}---a fundamental constraint on GNN improvement.

\textbf{Over-smoothing}: Li et al.~\cite{li2018oversmoothing} identified representation collapse in deep GNNs. Chen et al.~\cite{chen2020oversmoothing} analyzed the spectral perspective. Our ADR provides a related but distinct insight: even \textit{one} layer of aggregation can damage information at mid-homophily.

\subsection{GNN Expressiveness Theory}

\textbf{WL hierarchy}: Xu et al.~\cite{xu2019gin} proved GNNs are at most as expressive as 1-WL. Morris et al.~\cite{morris2019wl} extended this to higher-order WL. These results focus on \textit{distinguishing graphs}, not \textit{improving predictions}---a gap our framework addresses.

\textbf{Information-theoretic bounds}: Wu et al.~\cite{wu2023infognn} analyzed information flow in GNNs. Oono and Suzuki~\cite{oono2020exponential} proved exponential information loss in deep GNNs. Our Information Budget is complementary: it bounds improvement from the \textit{task} perspective, not the \textit{architecture} perspective.

\textbf{Aggregation analysis}: Cai and Wang~\cite{cai2020note} showed mean aggregation has limitations. Hamilton~\cite{hamilton2020book} compared aggregators extensively. Our Theorem~\ref{thm:concat} provides theoretical grounding for concatenation's superiority.

\subsection{Positioning of Our Work}

Our contribution differs from prior work in three ways:

\begin{enumerate}
    \item \textbf{Budget-first perspective}: Existing frameworks start from $h$ or structure. We start from MLP accuracy, establishing the ceiling before considering how to reach it.

    \item \textbf{Damage quantification}: ADR provides a scalar measure of aggregation harm, explaining GraphSAGE's robustness through feature preservation.

    \item \textbf{Heterophily diagnostic}: Rather than treating low-$h$ as a single regime, we measure whether label homophily recovers at two hops and test, rather than assume, how this statistic relates to trained models.
\end{enumerate}

Together, these form a unified decision framework that predicts GNN outcomes more accurately than homophily alone.
